% Source: https://artofproblemsolving.com/community/c3h1862524_geometry_section_11__problem_n11
% Stable identity: aops-topic-1862524-q01
\textbf{Construction.} Put $P$, $E$, $D$ on a line in this order (so $E$ is between $P$ and $D$), and put $A$ on ray $\overrightarrow{ED}$, beyond $D$. Choose $F$ on one side of line $PD$ and $C$ on the other side.

\begin{asy}
usepackage("amsmath");
size(7.5cm);
defaultpen(fontsize(10pt));
pair P=(0,0),Ep=(2,0),D=(3.6,0),A=(5,0),F=(3.4,1.9),C=(1.3,-1.7);
draw(P--(6.6,0),red+linewidth(0.8),EndArrow(6));
draw(P--(P+1.45*(C-P)),red+linewidth(0.8),EndArrow(6));
draw(Ep--(Ep+1.5*(F-Ep)),blue+linewidth(0.8),EndArrow(6));
draw(Ep--(6.6,0.0),deepgreen+linewidth(2.6));
dot(P);dot(Ep);dot(D);dot(A);dot(F);dot(C);
label("$P$",P,W);label("$E$",Ep,S);label("$D$",D,S);label("$A$",A,S);label("$F$",F,E);label("$C$",C,E);
label("$\angle DPC$",P+(1.25,-0.55),red);label("$\angle AEF$",Ep+(1.15,0.45),blue);
label("$\overrightarrow{ED}$",(5.8,0),N,deepgreen);
label("\mbox{}",(min(currentpicture,user=true)+max(currentpicture,user=true))/2,black); // keeps the text after the figure black when LaTeX typesets the labels inline
\end{asy}

Then $\angle AEF=\overrightarrow{EA}\cup\overrightarrow{EF}$ (blue, with $\overrightarrow{EA}=\overrightarrow{ED}$) and $\angle DPC=\overrightarrow{PD}\cup\overrightarrow{PC}$ (red).

\textbf{Proof that the intersection is $\overrightarrow{ED}$.} Intersect each ray of one angle with each ray of the other.
\begin{itemize}
\item $\overrightarrow{EA}\cap\overrightarrow{PD}=\overrightarrow{ED}$: $E$ lies on $\overrightarrow{PD}$, and $\overrightarrow{EA}=\overrightarrow{ED}$ points the same way, so $\overrightarrow{ED}\subseteq\overrightarrow{PD}$.
\item $\overrightarrow{EA}\cap\overrightarrow{PC}=\varnothing$: ray $\overrightarrow{PC}$ meets line $PD$ only at $P$, and $P$ is not on $\overrightarrow{ED}$ since $E$ is between $P$ and $D$.
\item $\overrightarrow{EF}\cap\overrightarrow{PD}=\{E\}$: ray $\overrightarrow{EF}$ meets line $PD$ only at $E$.
\item $\overrightarrow{EF}\cap\overrightarrow{PC}=\varnothing$: apart from its endpoint, $\overrightarrow{EF}$ lies on $F$'s side of line $PD$ and $\overrightarrow{PC}$ on $C$'s side. Their endpoints $E$ and $P$ are different points of line $PD$, and neither lies on the other ray.
\end{itemize}
The union of these four intersections is $\overrightarrow{ED}\cup\{E\}=\overrightarrow{ED}$, as required. The ray $\overrightarrow{ED}$ must be a side of $\angle AEF$, so $D$ has to lie on $\overrightarrow{EA}$ or on $\overrightarrow{EF}$; here we chose $\overrightarrow{EA}$. If $F$ and $C$ were on the same side of line $PD$, rays $\overrightarrow{EF}$ and $\overrightarrow{PC}$ could cross and add an extra point to the intersection.
